\documentclass[
  doc,
  floatsintext,
  longtable,
  a4paper,
  nolmodern,
  notxfonts,
  notimes,
  12pt,
  colorlinks=true,linkcolor=blue,citecolor=blue,urlcolor=blue]{apa7}

\usepackage{amsmath}
\usepackage{amssymb}

\geometry{inner=1in, outer=1in}
\fancyhfoffset[LE,RO]{0cm}


\usepackage[bidi=default]{babel}
\babelprovide[main,import]{spanish}


% get rid of language-specific shorthands (see #6817):
\let\LanguageShortHands\languageshorthands
\def\languageshorthands#1{}

\RequirePackage{longtable}
\RequirePackage{threeparttablex}

\makeatletter
\renewcommand{\paragraph}{\@startsection{paragraph}{4}{\parindent}%
	{0\baselineskip \@plus 0.2ex \@minus 0.2ex}%
	{-.5em}%
	{\normalfont\normalsize\bfseries\typesectitle}}

\renewcommand{\subparagraph}[1]{\@startsection{subparagraph}{5}{0.5em}%
	{0\baselineskip \@plus 0.2ex \@minus 0.2ex}%
	{-\z@\relax}%
	{\normalfont\normalsize\bfseries\itshape\hspace{\parindent}{#1}\textit{\addperi}}{\relax}}
\makeatother




\usepackage{longtable, booktabs, multirow, multicol, colortbl, hhline, caption, array, float, xpatch}
\usepackage{subcaption}


\renewcommand\thesubfigure{\Alph{subfigure}}
\setcounter{topnumber}{2}
\setcounter{bottomnumber}{2}
\setcounter{totalnumber}{4}
\renewcommand{\topfraction}{0.85}
\renewcommand{\bottomfraction}{0.85}
\renewcommand{\textfraction}{0.15}
\renewcommand{\floatpagefraction}{0.7}

\usepackage{tcolorbox}
\tcbuselibrary{listings,theorems, breakable, skins}
\usepackage{fontawesome5}

\definecolor{quarto-callout-color}{HTML}{909090}
\definecolor{quarto-callout-note-color}{HTML}{0758E5}
\definecolor{quarto-callout-important-color}{HTML}{CC1914}
\definecolor{quarto-callout-warning-color}{HTML}{EB9113}
\definecolor{quarto-callout-tip-color}{HTML}{00A047}
\definecolor{quarto-callout-caution-color}{HTML}{FC5300}
\definecolor{quarto-callout-color-frame}{HTML}{ACACAC}
\definecolor{quarto-callout-note-color-frame}{HTML}{4582EC}
\definecolor{quarto-callout-important-color-frame}{HTML}{D9534F}
\definecolor{quarto-callout-warning-color-frame}{HTML}{F0AD4E}
\definecolor{quarto-callout-tip-color-frame}{HTML}{02B875}
\definecolor{quarto-callout-caution-color-frame}{HTML}{FD7E14}

%\newlength\Oldarrayrulewidth
%\newlength\Oldtabcolsep


\usepackage{hyperref}



\usepackage{color}
\usepackage{fancyvrb}
\newcommand{\VerbBar}{|}
\newcommand{\VERB}{\Verb[commandchars=\\\{\}]}
\DefineVerbatimEnvironment{Highlighting}{Verbatim}{commandchars=\\\{\}}
% Add ',fontsize=\small' for more characters per line
\usepackage{framed}
\definecolor{shadecolor}{RGB}{241,243,245}
\newenvironment{Shaded}{\begin{snugshade}}{\end{snugshade}}
\newcommand{\AlertTok}[1]{\textcolor[rgb]{0.68,0.00,0.00}{#1}}
\newcommand{\AnnotationTok}[1]{\textcolor[rgb]{0.37,0.37,0.37}{#1}}
\newcommand{\AttributeTok}[1]{\textcolor[rgb]{0.40,0.45,0.13}{#1}}
\newcommand{\BaseNTok}[1]{\textcolor[rgb]{0.68,0.00,0.00}{#1}}
\newcommand{\BuiltInTok}[1]{\textcolor[rgb]{0.00,0.23,0.31}{#1}}
\newcommand{\CharTok}[1]{\textcolor[rgb]{0.13,0.47,0.30}{#1}}
\newcommand{\CommentTok}[1]{\textcolor[rgb]{0.37,0.37,0.37}{#1}}
\newcommand{\CommentVarTok}[1]{\textcolor[rgb]{0.37,0.37,0.37}{\textit{#1}}}
\newcommand{\ConstantTok}[1]{\textcolor[rgb]{0.56,0.35,0.01}{#1}}
\newcommand{\ControlFlowTok}[1]{\textcolor[rgb]{0.00,0.23,0.31}{\textbf{#1}}}
\newcommand{\DataTypeTok}[1]{\textcolor[rgb]{0.68,0.00,0.00}{#1}}
\newcommand{\DecValTok}[1]{\textcolor[rgb]{0.68,0.00,0.00}{#1}}
\newcommand{\DocumentationTok}[1]{\textcolor[rgb]{0.37,0.37,0.37}{\textit{#1}}}
\newcommand{\ErrorTok}[1]{\textcolor[rgb]{0.68,0.00,0.00}{#1}}
\newcommand{\ExtensionTok}[1]{\textcolor[rgb]{0.00,0.23,0.31}{#1}}
\newcommand{\FloatTok}[1]{\textcolor[rgb]{0.68,0.00,0.00}{#1}}
\newcommand{\FunctionTok}[1]{\textcolor[rgb]{0.28,0.35,0.67}{#1}}
\newcommand{\ImportTok}[1]{\textcolor[rgb]{0.00,0.46,0.62}{#1}}
\newcommand{\InformationTok}[1]{\textcolor[rgb]{0.37,0.37,0.37}{#1}}
\newcommand{\KeywordTok}[1]{\textcolor[rgb]{0.00,0.23,0.31}{\textbf{#1}}}
\newcommand{\NormalTok}[1]{\textcolor[rgb]{0.00,0.23,0.31}{#1}}
\newcommand{\OperatorTok}[1]{\textcolor[rgb]{0.37,0.37,0.37}{#1}}
\newcommand{\OtherTok}[1]{\textcolor[rgb]{0.00,0.23,0.31}{#1}}
\newcommand{\PreprocessorTok}[1]{\textcolor[rgb]{0.68,0.00,0.00}{#1}}
\newcommand{\RegionMarkerTok}[1]{\textcolor[rgb]{0.00,0.23,0.31}{#1}}
\newcommand{\SpecialCharTok}[1]{\textcolor[rgb]{0.37,0.37,0.37}{#1}}
\newcommand{\SpecialStringTok}[1]{\textcolor[rgb]{0.13,0.47,0.30}{#1}}
\newcommand{\StringTok}[1]{\textcolor[rgb]{0.13,0.47,0.30}{#1}}
\newcommand{\VariableTok}[1]{\textcolor[rgb]{0.07,0.07,0.07}{#1}}
\newcommand{\VerbatimStringTok}[1]{\textcolor[rgb]{0.13,0.47,0.30}{#1}}
\newcommand{\WarningTok}[1]{\textcolor[rgb]{0.37,0.37,0.37}{\textit{#1}}}

\providecommand{\tightlist}{%
  \setlength{\itemsep}{0pt}\setlength{\parskip}{0pt}}
\usepackage{longtable,booktabs,array}
\newcounter{none} % for unnumbered tables
\usepackage{calc} % for calculating minipage widths
% Correct order of tables after \paragraph or \subparagraph
\usepackage{etoolbox}
\makeatletter
\patchcmd\longtable{\par}{\if@noskipsec\mbox{}\fi\par}{}{}
\makeatother
% Allow footnotes in longtable head/foot
\IfFileExists{footnotehyper.sty}{\usepackage{footnotehyper}}{\usepackage{footnote}}
\makesavenoteenv{longtable}

\usepackage{graphicx}
\makeatletter
\newsavebox\pandoc@box
\newcommand*\pandocbounded[1]{% scales image to fit in text height/width
  \sbox\pandoc@box{#1}%
  \Gscale@div\@tempa{\textheight}{\dimexpr\ht\pandoc@box+\dp\pandoc@box\relax}%
  \Gscale@div\@tempb{\linewidth}{\wd\pandoc@box}%
  \ifdim\@tempb\p@<\@tempa\p@\let\@tempa\@tempb\fi% select the smaller of both
  \ifdim\@tempa\p@<\p@\scalebox{\@tempa}{\usebox\pandoc@box}%
  \else\usebox{\pandoc@box}%
  \fi%
}
% Set default figure placement to htbp
\def\fps@figure{htbp}
\makeatother







\usepackage{newtx}

\defaultfontfeatures{Scale=MatchLowercase}
\defaultfontfeatures[\rmfamily]{Ligatures=TeX,Scale=1}





\title{Números}


\shorttitle{NÚMEROS}


\usepackage{etoolbox}



\ccoppy{\textcopyright~2026}



\author{Edison Achalma}



\affiliation{
{Escuela Profesional de Economía, Universidad Nacional de San Cristóbal
de Huamanga}}




\leftheader{Achalma}

\date{2026-10-20}


\abstract{Este abstract será actualizado una vez que se complete el
contenido final del artículo. }



\authornote{\par{\addORCIDlink{Edison Achalma}{0000-0001-6996-3364}} 
\par{ }
\par{   El autor no tiene conflictos de interés que revelar.    Los
roles de autor se clasificaron utilizando la taxonomía de roles de
colaborador (CRediT; https://credit.niso.org/) de la siguiente
manera:  Edison Achalma:   conceptualización, metodología, análisis
formal, investigación, recursos, curación de
datos, redacción, visualización, supervisión, administración del
proyecto}
\par{La correspondencia relativa a este artículo debe dirigirse a Edison
Achalma, Escuela Profesional de Economía, Universidad Nacional de San
Cristóbal de Huamanga, Portal Independencia N°
57, Ayacucho, AYA 5001, Perú, Email: \href{mailto:elmer.achalma.09@unsch.edu.pe}{elmer.achalma.09@unsch.edu.pe}}
}

\makeatletter
\let\endoldlt\endlongtable
\def\endlongtable{
\hline
\endoldlt
}
\makeatother

\urlstyle{same}



\hypersetup{
  pdftitle={N\'{u}meros},
  pdfauthor={Autor}
}
\makeatletter
\@ifpackageloaded{caption}{}{\usepackage{caption}}
\AtBeginDocument{%
\ifdefined\contentsname
  \renewcommand*\contentsname{Tabla de contenidos}
\else
  \newcommand\contentsname{Tabla de contenidos}
\fi
\ifdefined\listfigurename
  \renewcommand*\listfigurename{Lista de Figuras}
\else
  \newcommand\listfigurename{Lista de Figuras}
\fi
\ifdefined\listtablename
  \renewcommand*\listtablename{Lista de Tablas}
\else
  \newcommand\listtablename{Lista de Tablas}
\fi
\ifdefined\figurename
  \renewcommand*\figurename{Figura}
\else
  \newcommand\figurename{Figura}
\fi
\ifdefined\tablename
  \renewcommand*\tablename{Tabla}
\else
  \newcommand\tablename{Tabla}
\fi
}
\@ifpackageloaded{float}{}{\usepackage{float}}
\floatstyle{ruled}
\@ifundefined{c@chapter}{\newfloat{codelisting}{h}{lop}}{\newfloat{codelisting}{h}{lop}[chapter]}
\floatname{codelisting}{Listado}
\newcommand*\listoflistings{\listof{codelisting}{Lista de Listados}}
\makeatother
\makeatletter
\makeatother
\makeatletter
\@ifpackageloaded{caption}{}{\usepackage{caption}}
\@ifpackageloaded{subcaption}{}{\usepackage{subcaption}}
\makeatother
\makeatletter
\@ifpackageloaded{fontawesome5}{}{\usepackage{fontawesome5}}
\makeatother

% From https://tex.stackexchange.com/a/645996/211326
%%% apa7 doesn't want to add appendix section titles in the toc
%%% let's make it do it
\makeatletter
\xpatchcmd{\appendix}
  {\par}
  {\addcontentsline{toc}{section}{\@currentlabelname}\par}
  {}{}
\makeatother

%% Disable longtable counter
%% https://tex.stackexchange.com/a/248395/211326

\usepackage{etoolbox}

\makeatletter
\patchcmd{\LT@caption}
  {\bgroup}
  {\bgroup\global\LTpatch@captiontrue}
  {}{}
\patchcmd{\longtable}
  {\par}
  {\par\global\LTpatch@captionfalse}
  {}{}
\apptocmd{\endlongtable}
  {\ifLTpatch@caption\else\addtocounter{table}{-1}\fi}
  {}{}
\newif\ifLTpatch@caption
\makeatother

\begin{document}

\maketitle


\hypertarget{toc}{}
\tableofcontents
\newpage
\section[Introduction]{Números}

\setcounter{secnumdepth}{3}

\setlength\LTleft{0pt}




\subsection{🎯 Motivación}\label{motivaciuxf3n}

Los números son fundamentales en programación: desde operaciones
aritméticas básicas hasta cálculos científicos y análisis de datos.
Python maneja tres tipos numéricos principales, cada uno con propósitos
específicos. Entender sus diferencias te permite elegir el tipo correcto
para cada situación y evitar errores de precisión o conversión.

\subsection{💡 Intuición}\label{intuiciuxf3n}

Imagina tres tipos de contenedores para números:

\begin{itemize}
\tightlist
\item
  \textbf{int (enteros)}: como una hoja de papel donde escribes números
  sin decimales: 5, -3, 42.
\item
  \textbf{float (decimales)}: como una calculadora científica que maneja
  decimales y notación exponencial: 3.14, -0.5, 1.5e10.
\item
  \textbf{complex (complejos)}: como un plano cartesiano con dos ejes
  (real e imaginario): 3+5j.
\end{itemize}

Cada uno tiene sus propias reglas y usos, y Python te permite convertir
entre ellos cuando es necesario.

\subsection{🧠 Modelo mental}\label{modelo-mental}

\begin{itemize}
\tightlist
\item
  \textbf{int}: precisión exacta, sin límite de tamaño (puede ser tan
  grande como la memoria lo permita).
\item
  \textbf{float}: precisión limitada (aproximación decimal, típicamente
  15-17 dígitos significativos).
\item
  \textbf{complex}: dos floats (parte real e imaginaria) combinados.
\item
  \textbf{Conversión}: posible entre int y float (puede perder
  precisión), pero no directamente desde complex.
\end{itemize}

\subsection{📐 Desarrollo}\label{desarrollo}

\subsubsection{1. Tipos Numéricos en
Python}\label{tipos-numuxe9ricos-en-python}

{\def\LTcaptype{none} % do not increment counter
\begin{longtable}[]{@{}
  >{\raggedright\arraybackslash}p{(\linewidth - 6\tabcolsep) * \real{0.0780}}
  >{\raggedright\arraybackslash}p{(\linewidth - 6\tabcolsep) * \real{0.1702}}
  >{\raggedright\arraybackslash}p{(\linewidth - 6\tabcolsep) * \real{0.2979}}
  >{\raggedright\arraybackslash}p{(\linewidth - 6\tabcolsep) * \real{0.4539}}@{}}
\toprule\noalign{}
\begin{minipage}[b]{\linewidth}\raggedright
Tipo
\end{minipage} & \begin{minipage}[b]{\linewidth}\raggedright
Ejemplo
\end{minipage} & \begin{minipage}[b]{\linewidth}\raggedright
Descripción
\end{minipage} & \begin{minipage}[b]{\linewidth}\raggedright
Cuándo usar
\end{minipage} \\
\midrule\noalign{}
\endhead
\bottomrule\noalign{}
\endlastfoot
\textbf{int} & \texttt{5}, \texttt{-3}, \texttt{1000000} & Números
enteros sin decimales & Conteos, índices, cantidades exactas \\
\textbf{float} & \texttt{3.14}, \texttt{-0.5}, \texttt{1.5e10} & Números
con decimales (precisión limitada) & Mediciones, cálculos científicos,
dinero (cuidado con precisión) \\
\textbf{complex} & \texttt{3+5j}, \texttt{1j}, \texttt{-5j} & Números
complejos (real + imaginario) & Matemáticas avanzadas, ingeniería,
física \\
\end{longtable}
}

\subsubsection{2. Enteros (int)}\label{enteros-int}

Los enteros son números sin decimales, sin límite de tamaño:

\begin{Shaded}
\begin{Highlighting}[]
\CommentTok{\# Enteros positivos y negativos}
\NormalTok{x }\OperatorTok{=} \DecValTok{42}
\NormalTok{y }\OperatorTok{=} \OperatorTok{{-}}\DecValTok{7}
\NormalTok{z }\OperatorTok{=} \DecValTok{12345678901234567890}  \CommentTok{\# Puede ser enorme}

\BuiltInTok{print}\NormalTok{(}\BuiltInTok{type}\NormalTok{(x))  }\CommentTok{\# \textless{}class \textquotesingle{}int\textquotesingle{}\textgreater{}}
\BuiltInTok{print}\NormalTok{(}\BuiltInTok{type}\NormalTok{(z))  }\CommentTok{\# \textless{}class \textquotesingle{}int\textquotesingle{}\textgreater{}}
\end{Highlighting}
\end{Shaded}

\subsubsection{3. Floats (decimales)}\label{floats-decimales}

Los floats representan números con decimales:

\begin{Shaded}
\begin{Highlighting}[]
\CommentTok{\# Floats básicos}
\NormalTok{x }\OperatorTok{=} \FloatTok{3.14}
\NormalTok{y }\OperatorTok{=} \OperatorTok{{-}}\FloatTok{0.5}
\NormalTok{z }\OperatorTok{=} \FloatTok{1.0}

\BuiltInTok{print}\NormalTok{(}\BuiltInTok{type}\NormalTok{(x))  }\CommentTok{\# \textless{}class \textquotesingle{}float\textquotesingle{}\textgreater{}}

\CommentTok{\# Notación científica}
\NormalTok{a }\OperatorTok{=} \FloatTok{35e3}   \CommentTok{\# 35 * 10\^{}3 = 35000.0}
\NormalTok{b }\OperatorTok{=} \FloatTok{12E4}   \CommentTok{\# 12 * 10\^{}4 = 120000.0}
\NormalTok{c }\OperatorTok{=} \OperatorTok{{-}}\FloatTok{87.7e100}  \CommentTok{\# Número muy pequeño (científico)}

\BuiltInTok{print}\NormalTok{(a)  }\CommentTok{\# 35000.0}
\end{Highlighting}
\end{Shaded}

\subsubsection{4. Números Complejos
(complex)}\label{nuxfameros-complejos-complex}

Los números complejos tienen parte real e imaginaria:

\begin{Shaded}
\begin{Highlighting}[]
\CommentTok{\# Números complejos}
\NormalTok{x }\OperatorTok{=} \DecValTok{3} \OperatorTok{+} \OtherTok{5j}
\NormalTok{y }\OperatorTok{=} \OtherTok{5j}          \CommentTok{\# Solo parte imaginaria}
\NormalTok{z }\OperatorTok{=} \OperatorTok{{-}}\OtherTok{5j}         \CommentTok{\# Imaginario negativo}

\BuiltInTok{print}\NormalTok{(}\BuiltInTok{type}\NormalTok{(x))  }\CommentTok{\# \textless{}class \textquotesingle{}complex\textquotesingle{}\textgreater{}}
\BuiltInTok{print}\NormalTok{(x)        }\CommentTok{\# (3+5j)}
\BuiltInTok{print}\NormalTok{(x.real)   }\CommentTok{\# 3.0 (parte real)}
\BuiltInTok{print}\NormalTok{(x.imag)   }\CommentTok{\# 5.0 (parte imaginaria)}
\end{Highlighting}
\end{Shaded}

\subsubsection{5. Conversión entre Tipos
(Casting)}\label{conversiuxf3n-entre-tipos-casting}

Puedes convertir entre tipos numéricos:

\begin{Shaded}
\begin{Highlighting}[]
\CommentTok{\# Conversiones básicas}
\NormalTok{x }\OperatorTok{=} \DecValTok{5}        \CommentTok{\# int}
\NormalTok{y }\OperatorTok{=} \FloatTok{3.14}     \CommentTok{\# float}
\NormalTok{z }\OperatorTok{=} \OtherTok{1j}       \CommentTok{\# complex}

\CommentTok{\# int → float}
\NormalTok{a }\OperatorTok{=} \BuiltInTok{float}\NormalTok{(x)      }\CommentTok{\# 5.0}
\BuiltInTok{print}\NormalTok{(}\BuiltInTok{type}\NormalTok{(a))    }\CommentTok{\# \textless{}class \textquotesingle{}float\textquotesingle{}\textgreater{}}

\CommentTok{\# float → int (trunca decimales)}
\NormalTok{b }\OperatorTok{=} \BuiltInTok{int}\NormalTok{(y)        }\CommentTok{\# 3 (pierde .14)}
\BuiltInTok{print}\NormalTok{(}\BuiltInTok{type}\NormalTok{(b))    }\CommentTok{\# \textless{}class \textquotesingle{}int\textquotesingle{}\textgreater{}}

\CommentTok{\# int → complex}
\NormalTok{c }\OperatorTok{=} \BuiltInTok{complex}\NormalTok{(x)    }\CommentTok{\# (5+0j)}
\BuiltInTok{print}\NormalTok{(}\BuiltInTok{type}\NormalTok{(c))    }\CommentTok{\# \textless{}class \textquotesingle{}complex\textquotesingle{}\textgreater{}}

\CommentTok{\# ❌ No se puede convertir complex → int/float}
\CommentTok{\# d = int(z)      \# TypeError: can\textquotesingle{}t convert complex to int}
\end{Highlighting}
\end{Shaded}

\subsubsection{6. Operaciones Numéricas
Comunes}\label{operaciones-numuxe9ricas-comunes}

\begin{Shaded}
\begin{Highlighting}[]
\NormalTok{a }\OperatorTok{=} \DecValTok{10}
\NormalTok{b }\OperatorTok{=} \DecValTok{3}

\CommentTok{\# Operaciones básicas}
\BuiltInTok{print}\NormalTok{(a }\OperatorTok{+}\NormalTok{ b)   }\CommentTok{\# 13 (suma)}
\BuiltInTok{print}\NormalTok{(a }\OperatorTok{{-}}\NormalTok{ b)   }\CommentTok{\# 7  (resta)}
\BuiltInTok{print}\NormalTok{(a }\OperatorTok{*}\NormalTok{ b)   }\CommentTok{\# 30 (multiplicación)}
\BuiltInTok{print}\NormalTok{(a }\OperatorTok{/}\NormalTok{ b)   }\CommentTok{\# 3.333... (float)}
\BuiltInTok{print}\NormalTok{(a }\OperatorTok{//}\NormalTok{ b)  }\CommentTok{\# 3  (división entera)}
\BuiltInTok{print}\NormalTok{(a }\OperatorTok{\%}\NormalTok{ b)   }\CommentTok{\# 1  (resto/modulo)}
\BuiltInTok{print}\NormalTok{(a }\OperatorTok{**}\NormalTok{ b)  }\CommentTok{\# 1000 (potencia)}
\end{Highlighting}
\end{Shaded}

\subsubsection{7. Números Aleatorios}\label{nuxfameros-aleatorios}

Python tiene un módulo \texttt{random} para generar números aleatorios:

\begin{Shaded}
\begin{Highlighting}[]
\ImportTok{import}\NormalTok{ random}

\CommentTok{\# Número aleatorio entre 1 y 9 (inclusive)}
\NormalTok{numero }\OperatorTok{=}\NormalTok{ random.randrange(}\DecValTok{1}\NormalTok{, }\DecValTok{10}\NormalTok{)}
\BuiltInTok{print}\NormalTok{(numero)}

\CommentTok{\# Número aleatorio entre 0.0 y 1.0}
\BuiltInTok{print}\NormalTok{(random.random())}

\CommentTok{\# Número aleatorio entre 1 y 10 (float)}
\BuiltInTok{print}\NormalTok{(random.uniform(}\DecValTok{1}\NormalTok{, }\DecValTok{10}\NormalTok{))}
\end{Highlighting}
\end{Shaded}

\subsubsection{8. Precisión de Floats}\label{precisiuxf3n-de-floats}

⚠️ \textbf{Importante:} los floats tienen precisión limitada:

\begin{Shaded}
\begin{Highlighting}[]
\CommentTok{\# Problemas de precisión}
\BuiltInTok{print}\NormalTok{(}\FloatTok{0.1} \OperatorTok{+} \FloatTok{0.2}\NormalTok{)  }\CommentTok{\# 0.30000000000000004 (no es exactamente 0.3)}

\CommentTok{\# Comparar floats}
\CommentTok{\# ❌ No recomendado}
\ControlFlowTok{if} \FloatTok{0.1} \OperatorTok{+} \FloatTok{0.2} \OperatorTok{==} \FloatTok{0.3}\NormalTok{:}
    \BuiltInTok{print}\NormalTok{(}\StringTok{"Igual"}\NormalTok{)  }\CommentTok{\# No se ejecuta}

\CommentTok{\# ✅ Mejor usar tolerancia}
\NormalTok{tolerancia }\OperatorTok{=} \FloatTok{1e{-}10}
\ControlFlowTok{if} \BuiltInTok{abs}\NormalTok{((}\FloatTok{0.1} \OperatorTok{+} \FloatTok{0.2}\NormalTok{) }\OperatorTok{{-}} \FloatTok{0.3}\NormalTok{) }\OperatorTok{\textless{}}\NormalTok{ tolerancia:}
    \BuiltInTok{print}\NormalTok{(}\StringTok{"Aproximadamente igual"}\NormalTok{)}
\end{Highlighting}
\end{Shaded}

\subsection{🔬 Ejemplos}\label{ejemplos}

\textbf{Mínimo ---} Tipos numéricos básicos:

\begin{Shaded}
\begin{Highlighting}[]
\NormalTok{entero }\OperatorTok{=} \DecValTok{42}
\NormalTok{decimal }\OperatorTok{=} \FloatTok{3.14}
\NormalTok{complejo }\OperatorTok{=} \DecValTok{2} \OperatorTok{+} \OtherTok{3j}

\BuiltInTok{print}\NormalTok{(}\BuiltInTok{type}\NormalTok{(entero))   }\CommentTok{\# \textless{}class \textquotesingle{}int\textquotesingle{}\textgreater{}}
\BuiltInTok{print}\NormalTok{(}\BuiltInTok{type}\NormalTok{(decimal))  }\CommentTok{\# \textless{}class \textquotesingle{}float\textquotesingle{}\textgreater{}}
\BuiltInTok{print}\NormalTok{(}\BuiltInTok{type}\NormalTok{(complejo)) }\CommentTok{\# \textless{}class \textquotesingle{}complex\textquotesingle{}\textgreater{}}
\end{Highlighting}
\end{Shaded}

\textbf{Progresivo ---} Conversiones y operaciones:

\begin{Shaded}
\begin{Highlighting}[]
\NormalTok{precio }\OperatorTok{=} \FloatTok{19.99}
\NormalTok{cantidad }\OperatorTok{=} \DecValTok{3}
\NormalTok{total }\OperatorTok{=}\NormalTok{ precio }\OperatorTok{*}\NormalTok{ cantidad}
\BuiltInTok{print}\NormalTok{(total)  }\CommentTok{\# 59.97 (float)}

\CommentTok{\# Convertir a entero (trunca)}
\NormalTok{total\_entero }\OperatorTok{=} \BuiltInTok{int}\NormalTok{(total)}
\BuiltInTok{print}\NormalTok{(total\_entero)  }\CommentTok{\# 59}

\CommentTok{\# Redondear}
\NormalTok{total\_redondeado }\OperatorTok{=} \BuiltInTok{round}\NormalTok{(total)}
\BuiltInTok{print}\NormalTok{(total\_redondeado)  }\CommentTok{\# 60}
\end{Highlighting}
\end{Shaded}

\textbf{Real ---} Cálculo de IMC (Índice de Masa Corporal):

\begin{Shaded}
\begin{Highlighting}[]
\CommentTok{\# Datos del usuario}
\NormalTok{peso\_kg }\OperatorTok{=} \FloatTok{72.5}    \CommentTok{\# float}
\NormalTok{altura\_m }\OperatorTok{=} \FloatTok{1.75}   \CommentTok{\# float}

\CommentTok{\# Cálculo del IMC}
\NormalTok{imc }\OperatorTok{=}\NormalTok{ peso\_kg }\OperatorTok{/}\NormalTok{ (altura\_m }\OperatorTok{**} \DecValTok{2}\NormalTok{)}
\BuiltInTok{print}\NormalTok{(}\SpecialStringTok{f"IMC: }\SpecialCharTok{\{}\NormalTok{imc}\SpecialCharTok{:.2f\}}\SpecialStringTok{"}\NormalTok{)  }\CommentTok{\# 23.67 (con 2 decimales)}

\CommentTok{\# Categorización}
\ControlFlowTok{if}\NormalTok{ imc }\OperatorTok{\textless{}} \FloatTok{18.5}\NormalTok{:}
\NormalTok{    categoria }\OperatorTok{=} \StringTok{"Bajo peso"}
\ControlFlowTok{elif}\NormalTok{ imc }\OperatorTok{\textless{}} \DecValTok{25}\NormalTok{:}
\NormalTok{    categoria }\OperatorTok{=} \StringTok{"Peso normal"}
\ControlFlowTok{elif}\NormalTok{ imc }\OperatorTok{\textless{}} \DecValTok{30}\NormalTok{:}
\NormalTok{    categoria }\OperatorTok{=} \StringTok{"Sobrepeso"}
\ControlFlowTok{else}\NormalTok{:}
\NormalTok{    categoria }\OperatorTok{=} \StringTok{"Obesidad"}

\BuiltInTok{print}\NormalTok{(}\SpecialStringTok{f"Categoría: }\SpecialCharTok{\{}\NormalTok{categoria}\SpecialCharTok{\}}\SpecialStringTok{"}\NormalTok{)}
\end{Highlighting}
\end{Shaded}

\subsection{⚠️ Errores comunes}\label{errores-comunes}

{\def\LTcaptype{none} % do not increment counter
\begin{longtable}[]{@{}
  >{\raggedright\arraybackslash}p{(\linewidth - 4\tabcolsep) * \real{0.2586}}
  >{\raggedright\arraybackslash}p{(\linewidth - 4\tabcolsep) * \real{0.3103}}
  >{\raggedright\arraybackslash}p{(\linewidth - 4\tabcolsep) * \real{0.4310}}@{}}
\toprule\noalign{}
\begin{minipage}[b]{\linewidth}\raggedright
Error
\end{minipage} & \begin{minipage}[b]{\linewidth}\raggedright
Síntoma
\end{minipage} & \begin{minipage}[b]{\linewidth}\raggedright
Solución
\end{minipage} \\
\midrule\noalign{}
\endhead
\bottomrule\noalign{}
\endlastfoot
\textbf{Dividir por cero} & \texttt{ZeroDivisionError} & Verifica que el
denominador no sea 0 \\
\textbf{Precisión de floats} & Resultados inesperados
(\texttt{0.1\ +\ 0.2}) & Usa \texttt{Decimal} para dinero o compara con
tolerancia \\
\textbf{Conversión de complex} & \texttt{TypeError} & No se puede
convertir complex a int/float \\
\textbf{Concatenar número + string} & \texttt{TypeError} & Convierte:
\texttt{str(numero)} o usa f-strings \\
\textbf{División con ints} & \texttt{5\ /\ 2} da \texttt{2.5} (float),
no \texttt{2} & Usa \texttt{//} para división entera \\
\end{longtable}
}

\subsection{✅ Buenas prácticas}\label{buenas-pruxe1cticas}

\begin{enumerate}
\def\labelenumi{\arabic{enumi}.}
\item
  \textbf{Usa \texttt{int} para conteos e índices}, \texttt{float} para
  mediciones y cálculos.
\item
  \textbf{Para dinero}, considera \texttt{decimal.Decimal} en lugar de
  float:

\begin{Shaded}
\begin{Highlighting}[]
\ImportTok{from}\NormalTok{ decimal }\ImportTok{import}\NormalTok{ Decimal}
\NormalTok{precio }\OperatorTok{=}\NormalTok{ Decimal(}\StringTok{\textquotesingle{}19.99\textquotesingle{}}\NormalTok{)}
\end{Highlighting}
\end{Shaded}
\item
  \textbf{Usa \texttt{//} para división entera} y \texttt{\%} para
  obtener el resto:

\begin{Shaded}
\begin{Highlighting}[]
\NormalTok{minutos }\OperatorTok{=} \DecValTok{125}
\NormalTok{horas }\OperatorTok{=}\NormalTok{ minutos }\OperatorTok{//} \DecValTok{60}   \CommentTok{\# 2}
\NormalTok{minutos\_restantes }\OperatorTok{=}\NormalTok{ minutos }\OperatorTok{\%} \DecValTok{60}  \CommentTok{\# 5}
\end{Highlighting}
\end{Shaded}
\item
  \textbf{Redondea explícitamente} cuando necesites controlar la
  precisión:

\begin{Shaded}
\begin{Highlighting}[]
\NormalTok{resultado }\OperatorTok{=} \FloatTok{3.14159265}
\BuiltInTok{print}\NormalTok{(}\BuiltInTok{round}\NormalTok{(resultado, }\DecValTok{2}\NormalTok{))  }\CommentTok{\# 3.14}
\end{Highlighting}
\end{Shaded}
\item
  \textbf{Prefiere \texttt{isinstance()} para verificar tipos}:

\begin{Shaded}
\begin{Highlighting}[]
\ControlFlowTok{if} \BuiltInTok{isinstance}\NormalTok{(valor, (}\BuiltInTok{int}\NormalTok{, }\BuiltInTok{float}\NormalTok{)):}
    \BuiltInTok{print}\NormalTok{(}\StringTok{"Es numérico"}\NormalTok{)}
\end{Highlighting}
\end{Shaded}
\item
  \textbf{Usa \texttt{\_} para separar números grandes} (Python 3.6+):

\begin{Shaded}
\begin{Highlighting}[]
\NormalTok{poblacion }\OperatorTok{=} \DecValTok{1\_000\_000\_000}  \CommentTok{\# Más legible que 1000000000}
\end{Highlighting}
\end{Shaded}
\end{enumerate}

\subsection{📝 Ejercicio}\label{ejercicio}

\textbf{Intermedio} --- Calculadora de Prestación

Escribe un programa que:

\begin{enumerate}
\def\labelenumi{\arabic{enumi}.}
\tightlist
\item
  Pida al usuario (simula con variables):

  \begin{itemize}
  \tightlist
  \item
    Salario mensual (float)
  \item
    Años trabajados (int)
  \item
    Edad (int)
  \end{itemize}
\item
  Calcule:

  \begin{itemize}
  \tightlist
  \item
    Salario anual (salario * 12)
  \item
    Prestación por antigüedad: 5\% del salario anual por cada año
    trabajado
  \item
    Edad en 10 años
  \end{itemize}
\item
  Muestre todos los resultados con 2 decimales
\item
  Convierta el resultado de la prestación a entero para mostrar la
  cantidad ``redondeada''
\end{enumerate}

\begin{center}\rule{0.5\linewidth}{0.5pt}\end{center}

\textbf{Solución (no mires hasta haberlo intentado):}

\begin{Shaded}
\begin{Highlighting}[]
\CommentTok{\# Datos}
\NormalTok{salario\_mensual }\OperatorTok{=} \FloatTok{2500.50}
\NormalTok{años\_trabajados }\OperatorTok{=} \DecValTok{8}
\NormalTok{edad }\OperatorTok{=} \DecValTok{32}

\CommentTok{\# Cálculos}
\NormalTok{salario\_anual }\OperatorTok{=}\NormalTok{ salario\_mensual }\OperatorTok{*} \DecValTok{12}
\NormalTok{prestacion\_porcentaje }\OperatorTok{=} \FloatTok{0.05} \OperatorTok{*}\NormalTok{ años\_trabajados}
\NormalTok{prestacion }\OperatorTok{=}\NormalTok{ salario\_anual }\OperatorTok{*}\NormalTok{ prestacion\_porcentaje}
\NormalTok{edad\_futura }\OperatorTok{=}\NormalTok{ edad }\OperatorTok{+} \DecValTok{10}

\CommentTok{\# Mostrar resultados}
\BuiltInTok{print}\NormalTok{(}\StringTok{"=== Cálculo de Prestación ==="}\NormalTok{)}
\BuiltInTok{print}\NormalTok{(}\SpecialStringTok{f"Salario mensual: $}\SpecialCharTok{\{}\NormalTok{salario\_mensual}\SpecialCharTok{:.2f\}}\SpecialStringTok{"}\NormalTok{)}
\BuiltInTok{print}\NormalTok{(}\SpecialStringTok{f"Salario anual: $}\SpecialCharTok{\{}\NormalTok{salario\_anual}\SpecialCharTok{:.2f\}}\SpecialStringTok{"}\NormalTok{)}
\BuiltInTok{print}\NormalTok{(}\SpecialStringTok{f"Años trabajados: }\SpecialCharTok{\{}\NormalTok{a}\SpecialCharTok{ñ}\NormalTok{os\_trabajados}\SpecialCharTok{\}}\SpecialStringTok{"}\NormalTok{)}
\BuiltInTok{print}\NormalTok{(}\SpecialStringTok{f"Prestación: $}\SpecialCharTok{\{}\NormalTok{prestacion}\SpecialCharTok{:.2f\}}\SpecialStringTok{"}\NormalTok{)}
\BuiltInTok{print}\NormalTok{(}\SpecialStringTok{f"Prestación redondeada: $}\SpecialCharTok{\{}\BuiltInTok{int}\NormalTok{(prestacion)}\SpecialCharTok{\}}\SpecialStringTok{"}\NormalTok{)}
\BuiltInTok{print}\NormalTok{(}\SpecialStringTok{f"Edad actual: }\SpecialCharTok{\{}\NormalTok{edad}\SpecialCharTok{\}}\SpecialStringTok{"}\NormalTok{)}
\BuiltInTok{print}\NormalTok{(}\SpecialStringTok{f"Edad en 10 años: }\SpecialCharTok{\{}\NormalTok{edad\_futura}\SpecialCharTok{\}}\SpecialStringTok{"}\NormalTok{)}
\end{Highlighting}
\end{Shaded}

\subsection{🗂️ Síntesis}\label{suxedntesis}

\begin{itemize}
\tightlist
\item
  \textbf{\texttt{int}}: números enteros sin decimales, sin límite de
  tamaño.
\item
  \textbf{\texttt{float}}: números con decimales, precisión limitada.
\item
  \textbf{\texttt{complex}}: números complejos (real + imaginario, con
  \texttt{j}).
\item
  \textbf{\texttt{type()}}: verifica el tipo de un número.
\item
  \textbf{Conversiones}: \texttt{int()}, \texttt{float()},
  \texttt{complex()}.
\item
  \textbf{Operaciones}: \texttt{+}, \texttt{-}, \texttt{*}, \texttt{/},
  \texttt{//}, \texttt{\%}, \texttt{**}.
\item
  \textbf{Precisión}: los floats son aproximados (problemas con
  decimales).
\item
  \textbf{\texttt{random}}: módulo para generar números aleatorios.
\end{itemize}

\subsection{🔗 Qué sigue}\label{quuxe9-sigue}

\begin{itemize}
\tightlist
\item
  \textbf{Relacionado:} operadores matemáticos, módulo \texttt{math},
  formateo de números
\item
  \textbf{Siguiente tema:} Strings --- aprenderás a trabajar con texto y
  todas sus operaciones.
\end{itemize}

\section{Publicaciones Similares}\label{publicaciones-similares}

Si te interesó este artículo, te recomendamos que explores otros blogs y
recursos relacionados que pueden ampliar tus conocimientos. Aquí te dejo
algunas sugerencias:

\begin{enumerate}
\def\labelenumi{\arabic{enumi}.}
\tightlist
\item
  \href{https://numerus-scriptum.netlify.app/python/2020-06-19-instalacion-de-anaconda/index.pdf}{\faIcon{file-pdf}}
  \href{https://numerus-scriptum.netlify.app/python/2020-06-19-instalacion-de-anaconda}{Instalacion
  De Anaconda}
\item
  \href{https://numerus-scriptum.netlify.app/python/2020-06-20-configurar-entorno-virtual-python-anaconda/index.pdf}{\faIcon{file-pdf}}
  \href{https://numerus-scriptum.netlify.app/python/2020-06-20-configurar-entorno-virtual-python-anaconda}{Configurar
  Entorno Virtual Python Anaconda}
\item
  \href{https://numerus-scriptum.netlify.app/python/2021-04-17-01-introducion-a-la-programacion-con-python/index.pdf}{\faIcon{file-pdf}}
  \href{https://numerus-scriptum.netlify.app/python/2021-04-17-01-introducion-a-la-programacion-con-python}{01
  Introducion A La Programacion Con Python}
\item
  \href{https://numerus-scriptum.netlify.app/python/2021-05-31-02-variables-expresiones-y-statements-con-python/index.pdf}{\faIcon{file-pdf}}
  \href{https://numerus-scriptum.netlify.app/python/2021-05-31-02-variables-expresiones-y-statements-con-python}{02
  Variables Expresiones Y Statements Con Python}
\item
  \href{https://numerus-scriptum.netlify.app/python/2021-06-07-03-objetos-de-python/index.pdf}{\faIcon{file-pdf}}
  \href{https://numerus-scriptum.netlify.app/python/2021-06-07-03-objetos-de-python}{03
  Objetos De Python}
\item
  \href{https://numerus-scriptum.netlify.app/python/2021-06-14-04-ejecucion-condicional-con-python/index.pdf}{\faIcon{file-pdf}}
  \href{https://numerus-scriptum.netlify.app/python/2021-06-14-04-ejecucion-condicional-con-python}{04
  Ejecucion Condicional Con Python}
\item
  \href{https://numerus-scriptum.netlify.app/python/2021-06-21-05-iteraciones-con-python/index.pdf}{\faIcon{file-pdf}}
  \href{https://numerus-scriptum.netlify.app/python/2021-06-21-05-iteraciones-con-python}{05
  Iteraciones Con Python}
\item
  \href{https://numerus-scriptum.netlify.app/python/2021-08-16-06-funciones-con-python/index.pdf}{\faIcon{file-pdf}}
  \href{https://numerus-scriptum.netlify.app/python/2021-08-16-06-funciones-con-python}{06
  Funciones Con Python}
\item
  \href{https://numerus-scriptum.netlify.app/python/2021-08-23-07-dataframes-con-python/index.pdf}{\faIcon{file-pdf}}
  \href{https://numerus-scriptum.netlify.app/python/2021-08-23-07-dataframes-con-python}{07
  Dataframes Con Python}
\item
  \href{https://numerus-scriptum.netlify.app/python/2021-11-29-08-prediccion-y-metrica-de-performance-con-python/index.pdf}{\faIcon{file-pdf}}
  \href{https://numerus-scriptum.netlify.app/python/2021-11-29-08-prediccion-y-metrica-de-performance-con-python}{08
  Prediccion Y Metrica De Performance Con Python}
\item
  \href{https://numerus-scriptum.netlify.app/python/2021-12-06-09-metodos-de-machine-learning-para-clasificacion-con-python/index.pdf}{\faIcon{file-pdf}}
  \href{https://numerus-scriptum.netlify.app/python/2021-12-06-09-metodos-de-machine-learning-para-clasificacion-con-python}{09
  Metodos De Machine Learning Para Clasificacion Con Python}
\item
  \href{https://numerus-scriptum.netlify.app/python/2021-12-13-10-metodos-de-machine-learning-para-regresion-con-python/index.pdf}{\faIcon{file-pdf}}
  \href{https://numerus-scriptum.netlify.app/python/2021-12-13-10-metodos-de-machine-learning-para-regresion-con-python}{10
  Metodos De Machine Learning Para Regresion Con Python}
\item
  \href{https://numerus-scriptum.netlify.app/python/2025-05-10-visualizacion-de-datos-con-python/index.pdf}{\faIcon{file-pdf}}
  \href{https://numerus-scriptum.netlify.app/python/2025-05-10-visualizacion-de-datos-con-python}{Visualizacion
  De Datos Con Python}
\end{enumerate}

Esperamos que encuentres estas publicaciones igualmente interesantes y
útiles. ¡Disfruta de la lectura!






\end{document}
